Este projeto propõe medir em que proporção os kits de desenvolvimento de software (SDKs) de rastreamento de terceiros embarcados nos aplicativos Android governamentais brasileiros estão declarados nas políticas de privacidade e na seção ``Segurança dos dados'' da Google Play, tomando como unidade o par formado por um aplicativo e cada SDK nele encontrado. Com base em estudos que encontraram declarações ausentes ou incorretas em parte relevante dos aplicativos analisados, a hipótese é que ao menos 10\% dos pares aplicativo--SDK não terão declaração compatível. A pesquisa é quantitativa e observacional. A versão atual de cada aplicativo será baixada da Google Play e identificada pelo hash do arquivo; os SDKs serão identificados por três detectores, cuja concordância será medida pelo $\kappa$ de Fleiss; cada par aplicativo--SDK será classificado conforme o que os documentos declaram; e o tráfego de rede de uma subamostra será observado em emulador. A presença de um SDK no código não será tratada como prova de envio de dados nem de infração à LGPD. O estudo se limita ao Android, que concentra a maior parte dos celulares no Brasil e cujos pacotes podem ser obtidos e analisados com ferramentas abertas. Na preparação, foram identificados 396 candidatos governamentais, dos quais 376 receberam classificação preliminar de oficiais. Como o AndroZoo continha apenas 92 dos 396 candidatos, todos os aplicativos serão baixados da Google Play. No piloto com 20 pacotes, a proporção de aplicativos com rastreador variou de 75\% a 95\% conforme o detector. A revisão sistemática identificou 396 registros publicados de 2018 a 2026 com termos do tema no título, triou 312 após a remoção de duplicatas e encaminhou 148 à leitura completa, com 31 propostas de inclusão até o momento. A coleta principal e o teste da hipótese estão previstos para o Projeto Capstone II.
